\section{Introduction}
\label{sec:intro}

\IEEEPARstart{S}{olar} photovoltaic (PV) generation has grown rapidly across Latin America, with
Colombia reaching over 1~GW of installed capacity and targeting 9~GW by 2030
under its National Energy Plan \cite{upme2023}. Reliable short-term forecasts of
global horizontal irradiance (GHI) are a prerequisite for efficient dispatch of
PV-integrated power systems: they enable reserve scheduling, reduce imbalance
penalties, and support real-time grid management. In tropical countries such as
Colombia, the forecasting challenge is compounded by the high frequency and spatial
heterogeneity of convective cloud systems, which can reduce surface irradiance by
more than 80\% within minutes and are poorly captured by coarse-resolution numerical
weather prediction (NWP) models.

Satellite-based nowcasting methods bypass NWP limitations by inferring cloud
motion and optical depth directly from geostationary imagery. Instruments such as
the GOES-16 Advanced Baseline Imager (ABI) provide multi-spectral observations
every 10~minutes at sub-kilometre resolution, making them well-suited for
sub-hourly to hourly irradiance estimation and forecasting. The dominant paradigm
couples a spatial feature extractor (applied to a patch of satellite pixels
centred on the target site) with a temporal model that exploits the recent
history of observed GHI or satellite state \cite{zhang2015,reno2016}.

Deep learning has substantially advanced this paradigm. Convolutional neural
networks (CNNs) applied to satellite patches learn translation-invariant spatial
patterns associated with cloud regimes, and when combined with recurrent models
such as LSTMs they achieve state-of-the-art performance on benchmark datasets
\cite{alessandrini2015,huertas2018,paletta2021}. More recently, graph neural networks
(GNNs) have emerged as an alternative spatial encoder: by representing pixels as
nodes in a graph and learning neighbourhood aggregation rules, GNNs relax the
translation-invariance assumption of CNNs and can capture irregular spatial
dependencies introduced by non-uniform cloud structures \cite{hamilton2017,ji2022}.
However, systematic comparisons of GNN-based and CNN-based encoders for GHI
forecasting remain scarce, and none, to the best of our knowledge, address the
tropical Andean climate.

\subsection*{Contributions}

This paper makes the following contributions:

\begin{enumerate}
  \item We propose \textbf{GraphSAGE-LSTM}, a model that treats a GOES-16 satellite
        patch as an 8-connected pixel graph, applies GraphSAGE neighbourhood
        aggregation \cite{hamilton2017}, and passes the resulting sequence of
        graph embeddings through an LSTM to produce a point forecast.

  \item We conduct a \textbf{rigorous empirical comparison} across two sites,
        three forecast horizons (1, 3, and 6~hours), and four random seeds,
        including Bayesian hyperparameter optimisation via Optuna \cite{akiba2019}
        for both GraphSAGE-LSTM and ResNet-LSTM.

  \item We show that GraphSAGE-LSTM \textbf{consistently outperforms} ResNet-LSTM
        at the 6-hour horizon for the low-elevation high-irradiance site
        ($\skillday = 0.636$ vs.\ $0.614$), while the advantage diminishes for
        the high-altitude, high-variability site.

  \item We identify a \textbf{regime-dependence} of statistical vs.\ deep-learning
        methods: SARIMAX on the clear-sky index remains competitive at 6~hours
        for the tropical mountain site (Bogot{\'a}), where the regular diurnal
        cycle dominates over cloud-driven variability.

  \item We provide a \textbf{fully reproducible experimental framework} including
        manifests, patch extraction, model code, and evaluation scripts, which
        we release openly.
\end{enumerate}

\subsection*{Paper organisation}

Section~\ref{sec:related} reviews related work. Section~\ref{sec:methodology}
describes the datasets, model architectures, training protocol, and evaluation
metrics. Section~\ref{sec:results} presents quantitative results.
Section~\ref{sec:discussion} analyses key findings. Section~\ref{sec:conclusion}
concludes and outlines future work.
